% ============================================================
% ПРИЛОЖЕНИЯ — оформление по ДВФУ п. 5.1.6–5.1.7
% Каждое приложение начинается с новой страницы.
% Заголовок: «Приложение А» (по центру, полужирный),
% ниже — название приложения (по центру, полужирный).
% ============================================================
% \appendix

% Настройка нумерации: «Приложение А», «Приложение Б», ...
% \thesection содержит полный текст метки — он же идёт в оглавление
% \renewcommand{\thesection}{Приложение~\Asbuk{section}}

% Формат заголовка приложения по ДВФУ п. 5.1.6–5.1.7:
%   — первая строка: «Приложение А» (по центру, полужирный)
%   — вторая строка: название (по центру, полужирный)
%\titleformat{\section}[display]
%    {\filcenter\normalfont\normalsize\bfseries}
%    {\thesection}
%    {0pt}
%    {}
%\titlespacing*{\section}{0pt}{\baselineskip}{\baselineskip}

%\newpage\section{Исходный код ключевых модулей}
% Листинги: Stage, PipelineEngine, ModelRegistry, EventBus

%\newpage\section{Конфигурационные файлы}
% Пример config.yaml с маппингом язык → модель

%\newpage\section{Скриншоты интерфейса}
% Полноразмерные скриншоты всех экранов GUI

%\newpage\section{Результаты экспериментов}
% Полные таблицы WER/DER по всем конфигурациям

%\newpage\section{Акт внедрения / справка об использовании}
% (при наличии)
